\heading{Homework 2 · Part II}{Formative mock examination · Lectures 3 and 4}
\label{mock-start}
\textbf{Deadline: to be announced.}

\textbf{Time: 120 minutes. Diagnostic total: 100 points.}

After completing Part I, put it aside and attempt this part in one timed
sitting. Write by hand on paper. For an honest diagnostic, use only this
question paper, without notes, internet or AI; this is a practice recommendation,
not an announcement of the actual final exam's permitted materials.

The TA will mark this part seriously, but mistakes here do \textbf{not}
lower your final course grade. Keep the original attempt and label corrections
made after the timed sitting. Its score is separate from the assessed Part I score.

\begin{center}\begin{tabular}{lr}
Ten true/false claims, 5 points each & 50\\
Two goods, one welfare bill & 20\\
Keeping the emergency fund alive & 20\\
Two short answers, 5 points each & 10\\\hline
Total & 100
\end{tabular}\end{center}

\paragraph{True/false instructions.} State a verdict and justify it.
For a true claim give a valid argument; for a false one give a fully specified
counterexample or a rigorous contradiction. An unsupported label earns no credit.
Derivatives and optimisers must be evaluated in the relevant parameter regime.

\paragraph{Welfare convention.} For a price increase at fixed income m,
report CV, EV and the Marshallian consumer-surplus change as positive
\emph{losses}. Distinguish certain wealth from expected utility of random wealth.
